% !TEX root = ../main.tex
% Abstract --- drafted 2026-05-28 against
% manuscript/planning/abstract_drafting_spec.md
% (companion spec to conclusion_drafting_spec.md).
%
% Variant choices per slot:
%   Slot 1 (phenomenon):  phenomenon-first, BHL default
%   Slot 2 (tool):        PSALTer 'infeasible' framing,
%                         inline PGT introduction
%   Slot 3 (findings):    decoupling-first, sweep-second;
%                         chi_1 identification moved to Slot 4
%                         to remove duplication (Option B)
%   Slot 4 (closing):     'we find evidence' + dual verdict,
%                         structural-finding register
%
% Supervisor directions this session (binding):
%   - chi_1 named; delta_1 dropped (chi_1 is dominant in the
%     multi-coupling regime; delta_1 was key only in isolation)
%   - no magnitude qualifiers; plane-wave growth is tunable, so
%     physical magnitudes deferred to localised-geometry follow-up
%     per Discussion + Conclusion
%   - suppression NOT confined to parity-odd minimal coupling;
%     suppression is generic across the surveyed theory space
%
% Post-review revision pass 2026-05-28: tightened AI-register sentence
% tails, split the 51-w S2, fixed passive voice in S4/S5, replaced
% particle-physics "channel" with Gertsenshtein-corpus "conversion"/
% "effect", removed campaign-label "full sector-closure" per spec.
%
% Friend-review pass 2026-05-28: standard acronym order applied per
% PSALTer / PSALTer-v2 convention (full name in plain roman + italic
% \TIDAL{} macro in parens); PGT gloss inserted in S2a as a comma
% apposition, mirroring 2406.12826's "framework for modified gravity
% with torsion" framing.
\begin{abstract}
The Gertsenshtein effect --- the conversion of gravitational waves
into photons in a background magnetic field --- may provide one of
the few observational signatures of high-frequency gravitational
waves. This conversion has not been computed across the operator
space of Poincar\'e gauge theory (PGT), a framework for modified
gravity with torsion, and the size of that space puts a complete
manual derivation out of reach. We present `Tensor Integration and
Derivation for Any Lagrangian'~(\TIDAL{}), a symbolic--numerical
software package that derives and integrates the linearised equations
of motion of any tensorial Lagrangian on a given background. Torsion decouples from the
Gertsenshtein effect in the minimal-coupling limit, so the
conversion structure is carried entirely by the non-minimal couplings.
We survey the PGT action across the coupling
space, comparing each configuration against the Einstein--Maxwell
baseline. Both ghost-free amplification and suppression of the
conversion are discovered, and we
identify a $\MAGF{}\times\nabla\MAGT{}$ derivative coupling as the dominant
amplifying operator.
\end{abstract}
